\section{Limitations and conclusion}\label{sec:limitations}

\ribench{} covers three tasks rather than role integrity
exhaustively. Court and Parliament use fixed subordinate material
rather than interacting agents, and matched non-hierarchical controls
are absent, so the results do not isolate hierarchy as a cause.
\rev{OpenReview has two tiers but no exchange between them, so
clarification, feedback and iterative revision, which interactive
hierarchies depend on, are not tested.}
OpenReview's constructed-vote design jointly changes scores, stance,
and specificity\rev{; the paper-only, score-free, count-balanced and
ratio controls of \S\ref{sec:or} separate these, but their newly
generated reviews were checked mechanically rather than by blind rating}. Low Court error adoption does not distinguish record-based
verification from prior knowledge or establish robustness to errors
of legal reasoning. Parliament measures loss of both limitations only
among cited claims, not complete evidence preservation.

RI equally weights the avoidance of three specified failures, not
verified fulfillment of all obligations. Its interpretation depends
on the selected tasks, endpoints, and weights, and cross-scenario
differences may also reflect document length and task difficulty.
Court and Parliament scoring relies on one LLM judge with rule-based
fallback. Human validation uses three annotators
on stratified samples, so agreement is not a population accuracy
estimate, and the annotators themselves agree only slightly in Court and
fairly in Parliament. Capability comparisons involve few models, and the index entries are not matched to
our reasoning settings, so differing rankings do not show
independence from capability.

We introduce \ribench{} to evaluate role integrity through obligations
governing how senior agents handle subordinate material.
Across the evaluated models, senior agents rarely adopt a planted
procedural error but often follow generic review majorities or drop
witnesses' stated limitations. The tested obligation prompts mitigate
some failures without eliminating them and also change control
behavior and citation selection. These findings motivate direct,
obligation-specific evaluation of role integrity rather than
inferring it from general capability or explicit role instructions.
